\subsection{自由广群}

集合序列 \(\mathbf{M}_{n}(\mathbf{X})\) 按对整数 \(n \geqslant 1\) 的归纳定义如下：记 \(\mathbf{M}_{\mathbf{1}}(\mathbf{X}) = \mathbf{X}\) ，对于 \(n \geqslant 2\) ， \(\mathbf{M}_{n}(\mathbf{X})\) 是诸集合 \(\mathbf{M}_{p}(\mathbf{X}) \times \mathbf{M}_{n-p}(\mathbf{X})\) （ \(1 \leqslant p \leqslant n - 1\) ）之和。集合族 \((\mathbf{M}_{n}(\mathbf{X}))_{n \geqslant 1}\) 之和记为 \(\mathbf{M}(\mathbf{X})\) ；每个集合 \(\mathbf{M}_{n}(\mathbf{X})\) 都与其在 \(\mathbf{M}(\mathbf{X})\) 中的典范像等同。对 \(\mathbf{M}(\mathbf{X})\) 的每个元素 w，存在唯一的整数 n 使得 \(w \in \mathbf{M}_{n}(\mathbf{X})\) ；它称为 w 的长度，记为 \(l(w)\) 。集合 X 由 \(\mathbf{M}(\mathbf{X})\) 中长度为 1 的元素组成。

设 \(w\) 和 \(w'\) 属于 \(\mathbf{M}(\mathbf{X})\) ；记 \(p = l(w)\) ， \(q = l(w')\) 。 \((w, w')\) 在 \(\mathbf{M}_p(\mathbf{X}) \times \mathbf{M}_q(\mathbf{X})\) 到和集 \(\mathbf{M}_{p + q}(\mathbf{X})\) 的典范单射下的像称为 \(w\) 与 \(w'\) 的合成，记作 \(ww'\) 或 \(w.w'\) 。则 \(l(w.w') = l(w) + l(w')\) ，且 \(\mathbf{M}(\mathbf{X})\) 中每个长度 \(\geqslant 2\) 的元素都可以唯一地写成 \(w'w''\) 的形式，其中 \(w', w''\) 属于 \(\mathbf{M}(\mathbf{X})\) 。

集合 \(\mathbf{M}(\mathbf{X})\) 连同合成律 \((w, w') \mapsto w.w'\) 称为在 \(\mathbf{X}\) 上构造的自由广群（§1，第 1 号，定义 1）。

命题 1. 设 M 为一个广群。X 到 M 的每个映射 f 都可以以唯一的方式延拓为 M(X) 到 M 的一个态射。

对 \(n \geqslant\) 作归纳，映射 \(f_n: \mathbf{M}_n(\mathbf{X}) \to \mathbf{M}\) 定义如下：设 \(f_{1} = f\) ；对 \(n \geqslant 2\) ，映射 \(f_{n}\) 由 \(f_{n}(w, w') = f_{p}(w).f_{n-p}(w')\) 定义，其中 \(p = 1, 2, \ldots, n - 1\) ， \((w, w')\) 属于 \(\mathbf{M}_{p}(\mathbf{X}) \times \mathbf{M}_{n-p}(\mathbf{X})\) 。设 \(g\) 是 \(\mathbf{M}(\mathbf{X})\) 到 \(\mathbf{M}\) 的这样的映射：它诱导 \(f_{n}\) （在 \(\mathbf{M}_{n}(\mathbf{X})\) 上），对每个整数 \(n \geqslant 1\) 。显然， \(g\) 是 \(\mathbf{M}(\mathbf{X})\) 到 \(\mathbf{M}\) 的唯一态射，它延拓 \(f\) 。

设 \(u\) 是 \(\mathbf{X}\) 到一个集合 \(\mathbf{Y}\) 中的映射。由命题 1，存在唯一的 \(\mathbf{M}(\mathbf{X})\) 到 \(\mathbf{M}(\mathbf{Y})\) 中的同态，它与 \(u\) 在 \(\mathbf{X}\) 上一致。把它记作 \(\mathbf{M}(u)\) 。若 \(v\) 是 \(\mathbf{Y}\) 到一个集合 \(\mathbf{Z}\) 中的映射，则同态 \(\mathbf{M}(v) \circ \mathbf{M}(u)\) （ \(\mathbf{M}(\mathbf{X})\) 到 \(\mathbf{M}(\mathbf{Z})\) 中的）与 \(v \circ u\) 在 \(\mathbf{X}\) 上一致，从而

\[
\mathbf {M} (v) \circ \mathbf {M} (u) = \mathbf {M} (v \circ u).
\]

命题 2. 设 \(u: \mathbf{X} \to \mathbf{Y}\) 是一个映射。若 \(u\) 是单射（相应地，满射，双射），则 \(\mathbf{M}(u)\) 也是。

假设 \(u\) 是单射的。当 \(X\) 为空时， \(M(X)\) 为空，因此 \(M(u)\) 是单射的。若 \(X\) 非空，则存在一个映射 \(v\) （ \(Y\) 到 \(X\) 中的），使得 \(v \circ u\) 是 \(X\) 的恒等映射（集合论，II，§ 3，第 8 号，命题 8）；映射 \(M(v) \circ M(u) = M(v \circ u)\) 是 \(M(X)\) 的恒等映射，因此 \(M(u)\) 是单射的。

当 \(u\) 是满射时，存在一个映射 \(w\) （ \(Y\) 到 \(X\) 中的），使得 \(u \circ w\) 是 \(Y\) 的恒等映射（集合论，II，§3，第 8 号，命题 8）。于是 \(M(u) \circ M(w) = M(u \circ w)\) 是 \(M(Y)\) 的恒等映射，因此 \(M(u)\) 是满射的。

最后，如果 \(u\) 是双射的，则它是单射且满射的，因此 \(\mathbf{M}(u)\) 具有相同的性质。

设 S 为 X 的一个子集。根据命题 2，S 到 X 的包含映射可以延拓为 M(S) 到 M(X) 的一个子广群 M'(S) 上的同构。借助该同构，广群 M(S) 与 M'(S) 视为同一。于是 M(S) 就是由 S 生成的 M(X) 的子广群。

设 X 为一个集合， \((u_{\alpha}, v_{\alpha})_{\alpha \in I}\) 是 M(X) 的元素组成的有序对的一个族。设 R 是 M(X) 上与 M(X) 的合成律相容且由诸 \((u_{\alpha}, v_{\alpha})\) 生成的等价关系（§1，第 6 号）。广群 M(X)/R 称为由 X 和关系元 \((u_{\alpha}, v_{\alpha})_{\alpha \in I}\) 所定义的广群。设 h 是 M(X) 到 M(X)/R 上的典范态射。则 M(X)/R 由 \(h(X)\) 生成。

设 N 为一个广群， \((n_{x})_{x\in\mathbf{X}}\) 为 N 的元素的一个族。设 k 为从 \(\mathbf{M}(\mathbf{X})\) 到 N 的态射，使得 \(k(x)=n_{x}\) 对所有 \(x\in X\) 成立（命题 1）。若 \(k(u_{\alpha})=k(v_{\alpha})\) 对所有 \(\alpha\in I\) 成立，则存在唯一的态射 \(f\colon\mathbf{M}(\mathbf{X})/\mathbf{R}\to\mathbf{N}\) ，使得 \(f(h(x))=n_{x}\) 对所有 \(x\in X\) 成立（§ 1，第 6 号，命题 9）。

